\chapter{System Architecture, Data Engineering \& Interface Design}
\label{ch:system_architecture}

\section{End-to-End Architectural Topology}
EduTrack is engineered as a high-performance, modular academic management platform adhering to strict decoupling between the fundamental data structures, algorithmic engines, business domain services, and presentation interfaces. Figure~\ref{fig:system_layers} illustrates the five-tiered software architecture.

\begin{figure}[H]
\centering
\begin{tcolorbox}[colback=codebackground,colframe=primaryblue,title=\textbf{EduTrack Layered Architectural Topology}]
\begin{enumerate}[leftmargin=1.5cm]
    \item[\textbf{Tier 5:}] \textbf{Dual Presentation Layer}
    \begin{itemize}
        \item \texttt{EduTrackCLI}: Terminal-based interactive runner utilizing a standard numbered menu hierarchy ($0$--$8$) for headless server environments.
        \item \texttt{EduTrackGUI}: Desktop Swing application with an 8-tab visual workspace, customized card layouts, real-time tables, and status meters.
    \end{itemize}
    \item[\textbf{Tier 4:}] \textbf{Service Orchestration Layer}
    \begin{itemize}
        \item \texttt{AcademicSearchService}, \texttt{PlagiarismDetectionService}, \texttt{ResourceAllocationService}, \texttt{ExamSchedulingService}, \texttt{AnalyticsService}.
    \end{itemize}
    \item[\textbf{Tier 3:}] \textbf{Algorithmic Engine Core (Modules 1--6)}
    \begin{itemize}
        \item Strings (\texttt{KMP}, \texttt{Z-Algorithm}, \texttt{Rabin-Karp}, \texttt{Aho-Corasick}).
        \item Suffix Structures (\texttt{SuffixArray}, \texttt{SAIS}, \texttt{KasaiLCP}, \texttt{SuffixAutomaton}).
        \item Dynamic Programming (\texttt{Levenshtein}, \texttt{Damerau-Levenshtein}, \texttt{MCM}, \texttt{BitmaskDP}, \texttt{OBST}).
        \item Network Flow (\texttt{FordFulkerson}, \texttt{EdmondsKarp}, \texttt{DinicsAlgorithm}, \texttt{BipartiteMatching}, \texttt{KonigsTheorem}).
        \item NP-Completeness (\texttt{DpllSatSolver}, Reductions: $\text{SAT} \to \text{CLIQUE} \to \text{IS} \to \text{VC}$, \texttt{VertexCover2Approx}).
        \item Parallel \& Randomized (\texttt{RandomizedQuickSort}, \texttt{ReservoirSampling}, \texttt{MillerRabin}, \texttt{BlellochScan}, \texttt{ParallelReduce}).
    \end{itemize}
    \item[\textbf{Tier 2:}] \textbf{Domain Model \& Business Entities}
    \begin{itemize}
        \item \texttt{Course}, \texttt{Student}, \texttt{Faculty}, \texttt{AssignmentSubmission}, \texttt{Classroom}, \texttt{ActivityLog}.
    \end{itemize}
    \item[\textbf{Tier 1:}] \textbf{Foundational Data Structures (Zero \texttt{java.util.*})}
    \begin{itemize}
        \item \texttt{MyArrayList<T>}, \texttt{MyLinkedList<T>}, \texttt{MyQueue<T>}, \texttt{MyStack<T>}, \texttt{MyHashMap<K,V>}, \texttt{MyMinHeap<T>}, \texttt{DisjointSetUnion}, \texttt{Pair<A,B>}.
    \end{itemize}
\end{enumerate}
\end{tcolorbox}
\caption{Tiered Architectural Decomposition of the EduTrack Platform.}
\label{fig:system_layers}
\end{figure}

A foundational requirement is total encapsulation: no engine class, model entity, or algorithmic module is permitted to import or instantiate classes from \texttt{java.util.*} (with the sole exception of \texttt{java.util.Scanner} used strictly for keyboard input handling in the interactive CLI runner).

% ------------------------------------------------------------------------------
\section{Domain Model Specifications}
The domain layer encapsulates academic entities with complete internal collection tracking powered exclusively by \texttt{MyArrayList<T>} and \texttt{MyHashMap<K,V>}.

\subsection{Course Entity (\texttt{Course.java})}
The \texttt{Course} class represents university course catalog entries. It tracks course codes, titles, departments, credit weightings, maximum enrollment caps, and prerequisite dependencies. Table~\ref{tab:course_schema} details its attributes.

\begin{table}[H]
\centering
\small
\begin{tabularx}{\textwidth}{llllX}
\toprule
\textbf{Attribute} & \textbf{Data Type} & \textbf{Core Container} & \textbf{Nullability} & \textbf{Description} \\
\midrule
\texttt{code} & \texttt{String} & Primitive String & NOT NULL & Unique course identifier (e.g., \texttt{CS301}) \\
\texttt{title} & \texttt{String} & Primitive String & NOT NULL & Human-readable course title \\
\texttt{department} & \texttt{String} & Primitive String & NOT NULL & Academic offering department \\
\texttt{credits} & \texttt{int} & Primitive int & NOT NULL & Academic credit weight ($1$--$4$) \\
\texttt{semester} & \texttt{int} & Primitive int & NOT NULL & Recommended curricular semester \\
\texttt{capacity} & \texttt{int} & Primitive int & NOT NULL & Physical laboratory/classroom cap \\
\texttt{enrolledCount} & \texttt{int} & Primitive int & NOT NULL & Live count of registered students \\
\texttt{prerequisites} & \texttt{MyArrayList<String>} & Handcrafted Array & NOT NULL & Array of required course codes \\
\bottomrule
\end{tabularx}
\caption{Relational Schema for Course Entity.}
\label{tab:course_schema}
\end{table}

\subsection{Student Entity (\texttt{Student.java})}
The \texttt{Student} entity captures undergraduate and graduate profiles. It maintains cumulative grade point averages (CGPA), earned credits, and dynamic course registrations.

\begin{table}[H]
\centering
\small
\begin{tabularx}{\textwidth}{llllX}
\toprule
\textbf{Attribute} & \textbf{Data Type} & \textbf{Core Container} & \textbf{Nullability} & \textbf{Description} \\
\midrule
\texttt{studentId} & \texttt{String} & Primitive String & NOT NULL & University Roll Number (e.g., \texttt{S101}) \\
\texttt{name} & \texttt{String} & Primitive String & NOT NULL & Full legal student name \\
\texttt{email} & \texttt{String} & Primitive String & NOT NULL & Institutional email address \\
\texttt{department} & \texttt{String} & Primitive String & NOT NULL & Major department \\
\texttt{semester} & \texttt{int} & Primitive int & NOT NULL & Current academic semester \\
\texttt{gpa} & \texttt{double} & Primitive double & NOT NULL & Cumulative GPA on a 4.00 scale \\
\texttt{creditsEarned} & \texttt{int} & Primitive int & NOT NULL & Total verified credit units completed \\
\texttt{enrolledCourses} & \texttt{MyArrayList<String>} & Handcrafted Array & NOT NULL & Registered course code array \\
\bottomrule
\end{tabularx}
\caption{Relational Schema for Student Entity.}
\label{tab:student_schema}
\end{table}

\subsection{Faculty Entity (\texttt{Faculty.java})}
Faculty records specify academic ranks, departmental affiliations, domain expertise tags, and maximum semester course loads for allocation constraints.

\begin{table}[H]
\centering
\small
\begin{tabularx}{\textwidth}{llllX}
\toprule
\textbf{Attribute} & \textbf{Data Type} & \textbf{Core Container} & \textbf{Nullability} & \textbf{Description} \\
\midrule
\texttt{facultyId} & \texttt{String} & Primitive String & NOT NULL & Faculty employee ID (e.g., \texttt{F101}) \\
\texttt{name} & \texttt{String} & Primitive String & NOT NULL & Full academic name and title \\
\texttt{department} & \texttt{String} & Primitive String & NOT NULL & Parent academic department \\
\texttt{email} & \texttt{String} & Primitive String & NOT NULL & Official communication email \\
\texttt{expertiseAreas} & \texttt{MyArrayList<String>} & Handcrafted Array & NOT NULL & Array of research specialization keywords \\
\texttt{maxCourseLoad} & \texttt{int} & Primitive int & NOT NULL & Maximum allowed course teaching quota \\
\texttt{assignedCourses} & \texttt{MyArrayList<String>} & Handcrafted Array & NOT NULL & Live assigned course codes \\
\bottomrule
\end{tabularx}
\caption{Relational Schema for Faculty Entity.}
\label{tab:faculty_schema}
\end{table}

\subsection{Assignment Submission (\texttt{AssignmentSubmission.java})}
Maintains student document submissions, submission timestamps, full raw text, and cryptographic verification digests for academic integrity inspection.
\begin{itemize}
    \item \texttt{submissionId}: Unique submission ticket (e.g., \texttt{SUB-001}).
    \item \texttt{studentId}: Submitting student reference.
    \item \texttt{courseCode}: Course assignment context.
    \item \texttt{title}: Title of submitted project/essay.
    \item \texttt{submissionDate}: Timestamp string.
    \item \texttt{content}: Raw text corpus analyzed by string and suffix engines.
\end{itemize}

\subsection{Classroom Entity (\texttt{Classroom.java})}
Represents physical campus facilities including lecture halls and specialized computational laboratories, capturing room capacity, building codes, and hardware availability (\texttt{hasLabEquipment}).

\subsection{Activity Log Entity (\texttt{ActivityLog.java})}
Encapsulates real-time system audit events, capturing event identifiers, ISO timestamps, action classifications (e.g., \texttt{LOGIN}, \texttt{EXAM\_SUBMIT}, \texttt{GRADE\_CHANGE}), user foreign keys, and severity levels.

% ------------------------------------------------------------------------------
\section{Dataset Engineering \& Synthetic Benchmark Corpus}
EduTrack includes a comprehensive synthetic academic dataset located in \texttt{data/} structured to rigorously exercise all algorithmic paths.

\subsection{Courses Dataset (\texttt{data/courses.csv})}
The course catalog contains 10 realistic university courses spanning Computer Science, Mathematics, Electrical Engineering, and Humanities, incorporating prerequisite structures and capacity constraints (Table~\ref{tab:courses_data_sample}).

\begin{table}[H]
\centering
\small
\begin{tabularx}{\textwidth}{llcccl}
\toprule
\textbf{Code} & \textbf{Title} & \textbf{Dept} & \textbf{Credits} & \textbf{Capacity} & \textbf{Prerequisites} \\
\midrule
\texttt{CS101} & Introduction to Programming & CS & 4 & 60 & None \\
\texttt{CS201} & Data Structures and Algorithms & CS & 4 & 50 & \texttt{CS101} \\
\texttt{CS301} & Design \& Analysis of Algorithms & CS & 4 & 45 & \texttt{CS201} \\
\texttt{CS302} & Database Management Systems & CS & 3 & 45 & \texttt{CS201} \\
\texttt{CS401} & Artificial Intelligence & CS & 3 & 40 & \texttt{CS301} \\
\texttt{CS402} & Computer Networks & CS & 3 & 40 & \texttt{CS201} \\
\texttt{MA101} & Linear Algebra \& Calculus & MATH & 4 & 70 & None \\
\texttt{MA201} & Discrete Mathematics & MATH & 3 & 55 & \texttt{MA101} \\
\texttt{EE201} & Digital Logic \& Computer Org & EE & 4 & 50 & None \\
\texttt{HU101} & Professional Ethics \& Communication & HUM & 2 & 80 & None \\
\bottomrule
\end{tabularx}
\caption{Complete Course Catalog Corpus in EduTrack Dataset.}
\label{tab:courses_data_sample}
\end{table}

\subsection{Students Dataset (\texttt{data/students.csv})}
The student dataset comprises 12 distinct undergraduate profiles reflecting diverse GPAs ($3.40$--$3.98$), completed credits ($54$--$88$), and semester levels ($3$--$6$), enabling deterministic verification of Randomized QuickSort merit partitioning and reservoir sampling streams.

\subsection{Faculty Dataset (\texttt{data/faculty.csv})}
The faculty dataset contains 6 professors across CS, Mathematics, and EE. Each faculty record defines distinct domain expertise keywords (e.g., ``Algorithms'', ``Networks'', ``Databases'', ``Machine Learning'') and contractual course limits ($2$ to $3$ courses).

\subsection{Plagiarism Corpus (\texttt{data/assignments/})}
To benchmark multi-document plagiarism triage, three realistic essay submissions on \textit{``Ethical Implications of Autonomous Artificial Intelligence in Higher Education''} are provided:
\begin{enumerate}
    \item \textbf{Alice Smith (\texttt{alice\_ai_ethics.txt}, 1,180 chars):} Original academic treatise discussing algorithmic bias, automated grading accountability, and epistemic fairness.
    \item \textbf{Bob Jones (\texttt{bob\_ai_ethics.txt}, 1,215 chars):} Contains an identical 278-character verbatim excerpt copied directly from Alice's submission without citation:
    \begin{quote}
    \small\slshape
    ``Autonomous algorithms in educational assessment introduce systemic biases when training distributions fail to encapsulate neurodiverse response typologies. Fairness metrics must enforce statistical parity across demographic cohorts while guaranteeing individual calibration.''\dots
    \end{quote}
    \item \textbf{Carol White (\texttt{carol\_ai_ethics.txt}, 1,095 chars):} Negative control submission analyzing independent philosophical perspectives on AI ethics with no overlapping verbatim content.
\end{enumerate}

\subsection{Library Corpus (\texttt{data/library_docs/algorithms_handbook.txt})}
A multi-kilobyte academic text documenting fundamental concepts in algorithmic design, asymptotic notation, and network flow invariants. This corpus is indexed by the Suffix Array ($O(N)$ SA-IS) and Suffix Automaton (DAWG) to test sub-millisecond pattern lookup and LCP-based document clustering.

\subsection{Real-time Activity Stream (\texttt{data/activity_stream.log})}
Contains 14 streaming event records capturing timestamps, IP addresses, user identifiers, and operation codes. This stream drives the streaming Reservoir Sampling engine and Parallel Reduction tree diagnostics.

% ------------------------------------------------------------------------------
\section{Service Layer Orchestration}
The service layer coordinates domain entities and foundational algorithms to resolve complex university operational workflows.

\begin{figure}[H]
\centering
\begin{tcolorbox}[colback=codebackground,colframe=primaryblue,title=\textbf{EduTrack Service Layer Workflow Mapping}]
\begin{itemize}
    \item \textbf{\texttt{AcademicSearchService}:}
    \begin{itemize}
        \item Single pattern exact match $\to$ Knuth-Morris-Pratt (\texttt{KmpMatcher}).
        \item Repetitive prefix substring discovery $\to$ Z-Algorithm (\texttt{ZAlgorithm}).
        \item Course code and rolling token lookups $\to$ 64-bit Rabin-Karp (\texttt{RabinKarp}).
        \item Multi-keyword simultaneous triage $\to$ Aho-Corasick Automaton (\texttt{AhoCorasick}).
    \end{itemize}
    \item \textbf{\texttt{PlagiarismDetectionService}:}
    \begin{itemize}
        \item High-level cross-document similarity score $\to$ Levenshtein \& Damerau-Levenshtein Edit Distance.
        \item Exact contiguous verbatim plagiarism extraction $\to$ Generalized Suffix Array with Kasai LCP array.
        \item Instantaneous pattern verification against all submissions $\to$ Suffix Automaton (DAWG).
    \end{itemize}
    \item \textbf{\texttt{ResourceAllocationService}:}
    \begin{itemize}
        \item One-to-one faculty course assignment $\to$ Hopcroft-Karp Bipartite Matching.
        \item Multi-capacity lab workstation seat allocation $\to$ Dinic's Blocking Flow Network.
        \item Minimum bottleneck resource identification $\to$ König's Min-Cut Vertex Cover.
    \end{itemize}
    \item \textbf{\texttt{ExamSchedulingService}:}
    \begin{itemize}
        \item Timetable conflict satisfiability $\to$ DPLL 3-SAT Solver.
        \item Conflict reduction pipeline $\to$ $\text{3-SAT} \to \text{CLIQUE} \to \text{IS} \to \text{Vertex Cover}$.
        \item Minimal proctor assignment $\to$ 2-Approximation Vertex Cover via maximal matching.
    \end{itemize}
    \item \textbf{\texttt{AnalyticsService}:}
    \begin{itemize}
        \item Merit list ranking $\to$ Randomized QuickSort ($O(N \log N)$ expected).
        \item Live activity stream auditing $\to$ Reservoir Sampling (Algorithm R).
        \item High-throughput log metrics $\to$ Blelloch Parallel Prefix Scan \& Parallel Reduce.
        \item Asymmetric credential generation $\to$ Miller-Rabin Primality Verification.
    \end{itemize}
\end{itemize}
\end{tcolorbox}
\caption{Algorithmic Mapping across EduTrack Service Coordinators.}
\label{fig:service_mapping}
\end{figure}

% ------------------------------------------------------------------------------
\section{Dual Presentation Architecture: CLI \& GUI}

\subsection{Interactive Terminal CLI Runner (\texttt{EduTrackCLI.java})}
The terminal interface provides an administrative console for server environments. It features a numbered navigation menu:
\begin{itemize}
    \item \textbf{Option 1 (Academic Search Engine):} Executes KMP, Z-Algorithm, Rabin-Karp, and Aho-Corasick over course descriptions.
    \item \textbf{Option 2 (Plagiarism Detection):} Runs Levenshtein, Damerau-Levenshtein, and Suffix Array Kasai LCP over student essays.
    \item \textbf{Option 3 (Suffix Indexing):} Builds SA-IS and Suffix Automaton over the library handbook.
    \item \textbf{Option 4 (Resource Allocation):} Solves Bipartite Matching for faculty and Dinic's Flow for lab seat allocation.
    \item \textbf{Option 5 (Exam Timetabling):} Solves timetable conflict 3-SAT with DPLL and computes 2-approx proctor placement.
    \item \textbf{Option 6 (Parallel Analytics):} Executes Randomized QuickSort, Reservoir Sampling, and Blelloch Parallel Scan.
    \item \textbf{Option 7 (Verification Suite):} Executes the complete automated test harness validating all 16 algorithmic modules against regression assertions.
    \item \textbf{Option 8 (Algorithm Benchmark Arena):} Runs head-to-head algorithmic tournament showdowns with ASCII performance bars and speedup metrics.
    \item \textbf{Option 9 (Launch Desktop GUI):} Launches the graphical Swing application in a background process.
    \item \textbf{Option 0 (Exit):} Cleanly terminates system processes.
\end{itemize}

\subsection{Swing Desktop Graphical Interface (\texttt{EduTrackGUI.java})}
The graphical application provides a modern desktop interface built with Java Swing. Configured with a clean, high-contrast palette (Primary Blue \texttt{\#2563EB}, Dark Slate \texttt{\#0F172A}, Emerald Green \texttt{\#10B981}, and Off-White \texttt{\#F8FAFC}), the interface organizes functionality into 9 tabbed panels:
\begin{enumerate}
    \item \textbf{Tab 1: Academic Search Engine:} Interactive text inputs for course corpus queries with algorithmic selector (KMP, Z, Rabin-Karp, Aho-Corasick) and formatted match tables.
    \item \textbf{Tab 2: Plagiarism Detection Center:} Multi-document comparison matrix displaying Levenshtein distance, similarity percentage, and extracted verbatim text excerpts.
    \item \textbf{Tab 3: Suffix Structure Indexer:} Visual display of suffix arrays, LCP arrays, and DAWG state statistics for indexed academic texts.
    \item \textbf{Tab 4: Resource Allocation (Network Flow \& 2D Visualizer):} Embedded vector canvas rendering Source $S$, Faculty, Course, and Sink $T$ nodes with directed capacity/flow labels (`f/c`) and saturated edge highlights.
    \item \textbf{Tab 5: Exam Timetabling (NP-Completeness \& Graph Canvas):} 2D radial conflict graph visualizer with glowing gold proctor stations (\textbf{★}) verifying 2-approximation vertex covers.
    \item \textbf{Tab 6: Parallel Analytics \& Sorting:} Real-time merit ranking visualizer, Blelloch scan input/output comparison, and parallel reduction timing.
    \item \textbf{Tab 7: Cryptographic Security:} Key generation workbench with Miller-Rabin composite witness inspector and prime status indicators.
    \item \textbf{Tab 8: System Verification Suite:} One-click execution of the 16-algorithm test harness with live status badges (\texttt{[PASS]}), assertion outputs, and microsecond latencies.
    \item \textbf{Tab 9: Algorithm Benchmark Arena:} Interactive tournament ring featuring 5 showdown matchups (String Search, Network Flow, Suffix Indexing, DP Recurrences, Parallel/Sorting) with dynamic horizontal comparison bar charts and relative speedup badges (`⚡ 4.7x FASTER`).
\end{enumerate}

\subsection{Interactive 2D Visual Graph Engine (\texttt{GraphVisualizerPanel.java})}
To bridge abstract graph theory with intuitive visual verification, EduTrack features a handcrafted vector rendering engine built with Java \texttt{Graphics2D} (requiring zero third-party dependencies). It supports:
\begin{itemize}
    \item \textbf{Dynamic Flow Rendering:} Renders bipartite faculty assignments and Dinic layered networks. Directed edges display $(f / c)$ capacity indicators. Edges with active flow glow emerald green, while saturated bottleneck edges are highlighted in amber.
    \item \textbf{Vertex Cover Invariant Verification:} In exam conflict mode, all conflicting course pairs are linked by edges. Nodes selected by the 2-approximation algorithm receive glowing gold halo rings and a \texttt{"PROCTOR"} badge, visually demonstrating that every single conflict edge touches at least one covered vertex.
    \item \textbf{Interactive Geometry:} Supports mouse-hover incident edge highlighting, tooltips, and real-time layout recalculation.
\end{itemize}

\subsection{Benchmark Bar Chart Component (\texttt{BenchmarkBarChartPanel.java})}
A custom component that paints publication-quality horizontal comparative performance charts. It renders normalized execution bars, gradient fills, microsecond latency tags, complexity annotations, and automated speedup factor pills without relying on external charting libraries.
